We investigate the algebraic structure governing how errors compound through
multi-step mathematical reasoning. Modeling each proof step as a linear
perturbation $E_k = I + \eps A_k$ of the identity in $\GL_d(\R)$, we study
the Lyapunov exponent $\lambda$ of the random matrix product
$\prod_{k=1}^n E_k$ and the spectral gap of associated random walks on
finite groups.  Our main result establishes a Lyapunov exponent ordering:
when the error directions $\{A_k\}$ are simultaneously diagonalizable
(the abelian case), the Lyapunov exponent is strictly negative, whereas
non-simultaneously-diagonalizable (non-abelian) error directions produce
a strictly larger exponent, with the gap controlled at order $\Theta(\eps^2)$
by the expected squared Frobenius norm of the commutator $\E[\frobnorm{\comm{A_1}{A_2}}^2]$.
We complement this with an analysis of spectral gaps on Cayley graphs,
proving that the dihedral group $D_n$ with two generators has spectral gap
asymptotically four times that of the cyclic group $\Z_{2n}$ with matched
order.  Applied to iterative proof refinement, these results show that
commuting (independent) proof substeps minimize error accumulation, while
non-commuting operations can accelerate proof space exploration---two
fundamentally distinct optimization objectives.  All theoretical results
are verified by numerical experiments.
